\chapter{Introduction}
\label{chap:introduction}

This chapter introduces the central research problem, explains its academic and policy relevance, and defines the thesis scope and objectives. It places the study at the intersection of capital market liberalization and firm-level digital transformation, clarifies why MSCI A-share inclusion provides a credible quasi-experimental setting, and previews the empirical and methodological contributions developed in later chapters.

Beyond basic orientation, the chapter also states the analytical position used throughout the thesis: digital transformation is modeled as an investment process shaped by both internal managerial capability and external capital-market institutions. This perspective shifts the question from \emph{whether} firms transform to \emph{how} financial opening changes the conditions that make transformation feasible, credible, and persistent.

\section{Research Background}

The relationship between capital market liberalization and corporate digital transformation has become a central issue in development economics, corporate finance, and technology management. Over the last two decades, many emerging economies have gradually opened domestic capital markets to international investors. In parallel, firms have faced growing pressure to digitize production, operations, governance, and customer interfaces. Yet the two streams of research are still often studied in isolation. Liberalization studies usually focus on valuation, liquidity, and investment efficiency, while digital-transformation studies emphasize internal capabilities, leadership, and technology readiness~\cite{Fang2024}. The key question at their intersection remains open: can international capital-market integration raise firm-level digital transformation, and through which channels?

This thesis starts from a boundary question: digital transformation is not only a technological or managerial process, but also a financing and governance process shaped by market institutions. In many emerging markets, digital programs require high fixed costs, involve uncertain returns, unfold over long horizons, and depend on organizational learning complementarities~\cite{Chanias2019}. Even when expected social gains are high, these features can create financing frictions that delay or discourage adoption~\cite{ChenZhang2024}. Capital market liberalization may shift these investment conditions by broadening the investor base, lowering financing costs, and strengthening governance discipline~\cite{Barro1990}. Understanding this link is therefore important for both scholarship and policy.

The Chinese background is especially important for this research. China features scale, diversity, and institutional growth in ways that make it ideal for empirical study~\cite{Allen2017}. It is one of the world's largest stock markets and a leading hub for industrial digital transformation, yet capital account liberalization has mostly been gradual and driven by policy~\cite{PLi2023}. This slow process results in notable policy events with predictable timing. Among these, the inclusion of Chinese A-share firms in the MSCI Emerging Markets Index stands out as a key externally imposed reform that increased firms' exposure to international investors.

The practical motivation is equally strong. Policymakers increasingly treat digital transformation as a core source of productivity growth and competitiveness. However, firms differ substantially in their capacity to implement large-scale transformation programs~\cite{Liu2023}. If access to global capital and investor networks shapes adoption, market-opening reforms may generate broader technological spillovers than is usually assumed~\cite{Zheng2021}. At the same time, liberalization may widen regional and firm-level capability gaps if benefits accrue mainly to firms that already hold structural advantages. This tension makes careful empirical evaluation necessary.

Measurement is another key motivation. Much of the existing evidence relies on managerial surveys or narrow proxies such as IT spending and patent counts. These indicators often suffer from inconsistent coverage, self-reporting bias, and weak comparability across time~\cite{Wang2022}. To address this issue, the thesis constructs a firm-year digital-transformation index from annual-report disclosures using computational text analysis. Combining this high-coverage measurement strategy with a quasi-experimental design allows the study to provide evidence that is both identification-oriented and broad in empirical scope.

Chen et al.~\cite{Chen2022} show that increased analyst coverage and improved information accuracy are two mechanisms through which digital transformation strengthens firms' information environments. Cheng and Masron~\cite{ChengMasron2023} further argue that regulatory uncertainty can intensify firms' digital responses, highlighting the central role of external institutional conditions in shaping transformation behavior. Du and Wang~\cite{DuWang2024} find that the innovation benefits of digital transformation are significantly amplified by financial support, particularly through capital markets, with effects significant at the 1\% level. Zhang and Wang~\cite{ZhangWang2024} report that digitally transformed firms face lower equity financing costs, whereas Xue and Zhang~\cite{XueZhang2024} document improvements in debt-financing outcomes and overall firm performance. Taken together, this evidence suggests that the financial consequences of digital transformation are conditional on financing environment, disclosure quality, and governance structure.

This study builds on prior work while moving beyond correlation-based inference. Earlier studies often focus on one channel or one outcome and do not fully address treatment endogeneity. By contrast, this thesis combines mechanism-oriented analysis, a long panel (2010--2022), and a quasi-experimental policy shock to test whether liberalization is a causal driver of digital transformation intensity. Conceptually, digital transformation is treated as cumulative strategic investment: firms adjust governance, reassess financing constraints, and update expectations over time. Liberalization is expected to influence this path through at least three channels: financing-constraint relaxation, governance and disclosure upgrading through international oversight, and expanded access to cross-border knowledge networks. These channels anchor the empirical framework.

As a result, there are three main motivations. First, the thesis addresses a significant open question about the relationship between firm-level technological upgrading and global finance. Second, it provides methodologically rigorous evidence from a major emerging economy. Third, it offers policy-relevant insights for firms and regulators seeking to align capital-market reform with digital-development objectives.

\section{Research Questions and Objectives}

In light of the aforementioned motivation, this thesis addresses three relevant research problems.

\begin{itemize}
    \item First, does inclusion in the MSCI Emerging Markets Index, a proxy for capital market liberalization, significantly increase firm-level digital transformation intensity? This is the thesis's primary causal question. Specifically, the study evaluates whether treated firms exhibit a sustained post-inclusion increase in digital transformation relative to otherwise comparable untreated firms.
    \item Second, do treatment effects vary by technology domain, industry, region, and firm characteristics? Digital transformation is not a single technology, but a bundle of investments with distinct capital-intensity and capability requirements. The thesis tests whether effects are stronger in domains such as AI, big data, and cloud computing; whether responses differ across technology-intensive and traditional industries; and whether regional institutional quality shapes outcomes.
    \item Third, how does liberalization affect digital transformation? The thesis investigates three channels using theory and prior evidence: funding restrictions, governance and monitoring, and knowledge spillovers. In empirical implementation, this thesis directly tests the first two channels and treats the knowledge channel as theory-grounded but exploratory in the current sample.
\end{itemize}

Methodologically, the thesis aims to integrate validated NLP-based measurement with quasi-experimental econometrics in a transparent and reproducible large-panel framework. Substantively, it evaluates whether international financial integration can serve as a catalyst for firm-level technological upgrading. From a policy perspective, it provides evidence relevant to regulators, institutional investors, and corporate decision-makers seeking to understand the technology consequences of market-opening reforms.

The thesis follows five operational aims to achieve these objectives. It generates high-coverage firm-year digital transformation metrics from annual reports. It employs baseline and dynamic DID techniques to estimate treatment effects. It evaluates inferential stability through robustness and sensitivity tests. It uses heterogeneous effects to identify boundary conditions. It conducts mechanism evaluations to assess financing, governance, and knowledge channels.

\section{The Empirical Setting: China's A-Share Market and MSCI Index Inclusion}

The empirical background for this argument is the gradual inclusion of Chinese A-shares in the MSCI Emerging Markets Index, a watershed moment in China's financial opening process~\cite{YLiu2024}. Prior to 2018, global investors had limited access to China's domestic equity market due to ownership restrictions, quota systems, and market-access frictions. Although conduits such as QFII, RQFII, and the Stock Connect programs existed, broad passive global allocation to mainland A-shares remained limited~\cite{Awokuse2009}.

MSCI inclusion changed this environment by connecting eligible A-share companies to benchmark-driven international capital flows~\cite{Wei2020}. MSCI indexes are often used by global institutional investors and index-tracking funds. Being included in the MSCI Emerging Markets Index affects active portfolio building, analyst focus, and investor relations practices, in addition to passive investment flows~\cite{PLi2023}. The reform therefore modified firms' exposure to external financing and oversight conditions in a way that is economically significant.

The phased inclusion process is particularly useful for identification. Phase~I was implemented in June 2018 with 234 large-cap A-shares at a 5\% inclusion factor following the June 2017 announcement. By November 2019, Phase~II expanded coverage to 472 securities with a higher inclusion factor. This phased implementation provides clear treatment timing and observable pre- and post-policy windows. A key advantage of this setting is that inclusion criteria were largely rule-based and externally determined. Eligibility was based on market capitalization, liquidity, and Stock Connect accessibility rather than on discretionary judgments about firms' digital strategies~\cite{Yang2025}. Although no real-world setting is fully random, this institutional design strengthens the plausibility that inclusion captures exogenous variation in international capital access relative to firms' underlying digital transformation trajectories.

China is also an ideal context because firms operate under substantial heterogeneity in regional development, industry composition, ownership structures, and governance quality. This heterogeneity allows the thesis to test whether liberalization effects are uniform or conditional on absorptive capacity~\cite{Dong2023}. Zhang et al.~\cite{Zhang2024} show that treatment effects vary systematically with regional market development, ownership structure, and firm size, with stronger responses in certain non-Eastern and non-state-firm contexts. Embedding these moderation patterns within the MSCI setting helps move beyond average treatment effects toward a richer understanding of distributional outcomes.

\section{Overview of Research Design and Data}

The empirical design combines causal identification with large-scale computational measurement. The baseline model is a two-way fixed-effects DID that compares firms included in MSCI (treated) with firms not included (control) before and after inclusion. This baseline is complemented by event-study estimation for dynamic effects and pre-trend validation, propensity-score-matched DID for improved overlap on observables, and placebo-based falsification tests.

The data cover 35,264 firm-year observations from 2010 to 2022, representing 2,712 Chinese listed firms after sample cleaning. Core financial and governance variables are drawn from CSMAR, while treatment status is derived from MSCI constituent lists. Digital transformation outcomes are constructed from annual-report text using NLP-based keyword extraction across six technology domains: AI/machine learning, big data, cloud computing, IoT, blockchain/distributed ledger technologies, and enterprise digitization programs.

The baseline outcome is a composite index using log transformation of aggregated keyword frequencies. This approach addresses right skewness and enhances interpretability when disclosure dispersion is high. For robustness, alternative outcomes are also used, including domain-specific indices, normalized measures adjusted for report length, and additional transformations to assess sensitivity to scaling choices.

Mechanism analysis is embedded in the design rather than added as a descriptive task. Proxies related to financing constraints and governance quality are estimated through structured regressions to check whether channel patterns match the conceptual framework. The knowledge-spillover channel is retained as a theory-grounded mechanism and interpreted using exploratory evidence in the current sample.

Overall, the design is intentionally triangulated: one institutional shock, multiple estimators, multiple outcome definitions, and multiple mechanism checks. This triangulation increases confidence that the results reflect real economic effects rather than artifacts of any single model choice.

\section{Significance of Research}

This study is significant for three strands of literature: capital market liberalization, digital transformation economics, and computational empirical methods.

\begin{itemize}
    \item First, it expands the liberalization literature beyond traditional market results. Previous research often focuses on liquidity, valuation, return co-movement, or financing costs. This study demonstrates that liberalization can affect firms' technological paths, thus connecting financial integration to long-term productivity-related behavior.
    \item Second, it advances the digital transformation literature by highlighting external financial environments as a key factor in shaping corporate digital strategy. While much existing research focuses on internal capabilities, managerial cognition, and organizational change, this thesis demonstrates that access to international capital and investor oversight can significantly influence the speed and direction of transformation.
    \item Third, it offers mechanism-level evidence that directly tests financing and governance channels under a common treatment design, while treating knowledge spillovers as a theory-grounded exploratory channel with suggestive interpretation.
    \item Fourth, the thesis enhances methodology by combining quasi-experimental causal inference with validated Chinese-language text analytics at scale.
    \item Fifth, it contributes to the emerging-market comparative literature by offering high-resolution evidence from China.
    \item Sixth, it provides policy implications based on micro-level evidence. Regulators interested in the broader development impact of capital market reform can use these findings to assess whether market-opening policies promote strategic upgrading.
\end{itemize}

\section{Thesis Structure}

The thesis is organized in seven chapters, with each chapter designed to build logically on the previous one.

\begin{itemize}
    \item \textbf{Chapter~1} introduces the research problem, motivation, questions, objectives, empirical setting, and expected contributions.
    \item \textbf{Chapter~2} develops the literature review. It synthesizes foundational and recent studies on liberalization, digital transformation, mechanism channels, and text-based measurement.
    \item \textbf{Chapter~3} presents institutional background and formal hypothesis development. It details China's pre-liberalization capital market structure, the MSCI inclusion process, and the quasi-experimental properties of the reform.
    \item \textbf{Chapter~4} introduces Research Design~I focused on causal identification. It describes sample construction, variable definitions, treatment assignment, baseline and dynamic model specifications, and descriptive statistics.
    \item \textbf{Chapter~5} introduces Research Design~II focused on computational measurement. It documents corpus construction, Chinese text segmentation, dictionary design, index construction, and validity diagnostics.
    \item \textbf{Chapter~6} reports empirical results and discussion. It presents baseline effects, event-study dynamics, robustness checks, heterogeneity analyses, and mechanism evidence.
    \item \textbf{Chapter~7} concludes with a synthesis of findings, original contributions, limitations, and future research directions.
\end{itemize}
